\documentclass[11pt, a4paper, oneside]{article}
\usepackage[ngerman]{babel}
\usepackage{worksheet}

\begin{document}
    \author{L. Bung}
    \title{Ableitung von Polynomfunktionen \linebreak und Summenregel}
    \subject{Mathematik}
    \class{FHR2}

    \maketitle

    \task{Graphen der Ableitung von Polynomfunktionen}

    \begin{enumerate}[label=\alph*)]
        \item Zeichnen Sie jeweils den Funktionsgraphen der Ableitung ein.
        \item Bestimmen Sie zur eingezeichneten Ableitungsfunktion einen Funktionsterm.
    \end{enumerate}

    \pgfplotsset{
        ableitung/.style={
            axis lines=center, x=0.8cm, y=0.8cm, grid=major,
            xmin=-2.5, xmax=2.5, ymin=-3.5, ymax=3.5,
            xtick={-2,...,2}, ytick={-3,...,3},
            tick label style={font=\footnotesize},
            xlabel={$x$}, ylabel={$y$},
            xlabel style={above right}, ylabel style={above right},
            title style={yshift=2ex},
        }
    }

    \begin{center}
        \begin{tikzpicture}
            \begin{axis}[ableitung]
                \addplot[domain=-2.5:2.5, red, thick]{x};
                \node[red, below right] at (0,-.5) {$f(x)=x$};
            \end{axis}
        \end{tikzpicture}
        \hfill
        \begin{tikzpicture}
            \begin{axis}[ableitung]
                \addplot[domain=-1.87:1.87, red, thick, samples=100]{x^2};
                \node[red, right] at (-1.25,1.5) {$g(x)=x^2$};
            \end{axis}
        \end{tikzpicture}
        \hfill
        \begin{tikzpicture}
            \begin{axis}[ableitung]
                \addplot[domain=-1.52:1.52, red, thick, samples=100]{x^3};
                \node[red, right] at (-1,-1.5) {$h(x)=x^3$};
            \end{axis}
        \end{tikzpicture}
    \end{center}

    \checkered[5cm]

    \begin{hint}{Ableitung von Potenzfunktionen}
        Die Ableitung einer allgemeinen Potenzfunktion $f(x) = a \cdot x^n$ ist
        $f'(x) = a \cdot n \cdot x^{n-1}.$

        \vspace{12pt}
        \textbf{Beispiel}: Die Ableitung der Funktion $f(x) = 3x^2$ ist $f'(x) = 3 \cdot 2 \cdot x^{2-1} = 6x$.
    \end{hint}


    \pagebreak
    \task{Funktionsterme der Ableitung von Potenzfunktionen}

    \begin{enumerate}[label=\alph*)]
        \item Bestimmen Sie jeweils den Funktionsterm der Ableitung.
        \begin{multicols}{3}
            \begin{enumerate}[label=(\roman*)]
                \item $f(x)=x^4$
                \item $g(x)=3x^3$
                \item $h(x)=-\frac{1}{2}x^4$
            \end{enumerate}
        \end{multicols}

        \item Berechnen Sie die Steigung der Funktionen $f(x)$, $g(x)$ und $h(x)$ in den Punkten $x_1 = 2$, $x_2 = -1$ und $x_3=10$.
    \end{enumerate}

    \checkered[4.5cm]

    
    \begin{hint}{Summenregel}
        Für Funktionen der Form $f(x) = u(x) + v(x)$ gilt $f'(x) = u'(x) + v'(x)$.
    \end{hint}

    \task{Ableitung von Polynomfunktionen}

    \begin{enumerate}[label=\alph*)]
        \item Bestimmen Sie jeweils den Funktionsterm der Ableitung.
        \begin{enumerate}[label=(\roman*)]
            \item $f(x)=2x^2+x$
            \item $g(x)=-4x+\frac{1}{2}x^3$
            \item $h(x)=-\frac{1}{3}x^3+2x^2-\frac{5}{2}x-5$
        \end{enumerate}

        \item Berechnen Sie die Steigung der Funktionen $f(x)$, $g(x)$ und $h(x)$ in den Punkten $x_1 = 0$, $x_2 = -3$ und $x_3=3$.
    \end{enumerate}

    \checkered[4.5cm]

    \let\clearpage\relax
\end{document}
